% cap04/4.2_modelo_c4.tex
\section{Evoluci\'on del Modelado Topol\'ogico: El Est\'andar C4}

La conceptualizaci\'on arquitect\'onica de \textit{Democra} repudia categ\'oricamente el uso obsoleto y geom\'etricamente ambiguo de los tradicionales Diagramas de Arquitectura Generales o los cl\'asicos Diagramas Entidad-Relaci\'on (ER) para representar infraestructuras en la nube. La ingenier\'ia de software contempor\'anea exige una topolog\'ia matem\'atica y ruteada, desprovista de interpretaciones subjetivas. Por ello, este proyecto adopta el est\'andar de la industria conocido como el \textbf{Modelo C4} (\textit{Context, Containers, Components, and Code}), ideado por Simon Brown. Este modelo abstrae el software en niveles de \textit{zoom} geom\'etrico, garantizando que la documentaci\'on sea tan r\'igida y exacta como el c\'odigo fuente subyacente.

\subsection{Nivel 1: Contexto de Sistema (System Context)}

El primer estrato geom\'etrico de abstracci\'on se denomina el \textbf{Nivel 1: Contexto de Sistema}. En esta fase, se traza el per\'imetro absoluto del ecosistema de software. El \textit{Contexto de Sistema} representa a la plataforma \textit{Democra} como una caja negra de alto nivel, rodeada por actores humanos y sistemas satelitales externos. Este nivel erradica la ambig\"uedad al obligar al ingeniero de software a mapear expl\'icitamente cada interacci\'on externa, definiendo los l\'imites transaccionales del proyecto.

La utilidad de este diagrama de contexto es suprema: permite a los \textit{stakeholders} (sean directivos de ONGs o ingenieros de DevOps) visualizar de un vistazo qu\'e entidades de red interact\'uan con la aplicaci\'on. No se enfoca en tecnolog\'ias espec\'ificas (no se menciona a React ni a Node.js en este nivel), sino en la transferencia de valor.

\subsection{Clasificaci\'on de Actores (Actors) en el Ecosistema}

El Modelo C4 exige la identificaci\'on de los \textbf{Actores (Actors)}. En el caso de \textit{Democra}, el sistema de usuarios de la ONG (antes documentados gen\'ericamente) se clasifica de forma formal bajo roles de inyecci\'on transaccional:
\begin{itemize}
    \item \textbf{Administrador de Tenant (Actor Privilegiado):} Representa al gerente de la ONG, quien orquesta las operaciones, emite directrices y extrae auditor\'ias.
    \item \textbf{Voluntario de Campo (Actor Operativo):} Representa al usuario m\'ovil, que inyecta datos de recolecci\'on (ej. firmas, donaciones) en tiempo real desde zonas de baja conectividad.
\end{itemize}

Cada actor est\'a vinculado al Contexto de Sistema mediante l\'ineas unidireccionales de dependencia, documentando expl\'icitamente la intenci\'on del usuario y el protocolo de conexi\'on (ej. \textit{HTTPS/TLS}).

\subsection{Mapeo de Entidades Satelitales y Dependencias Externas}

El rigor del Modelo C4 proh\'ibe la omisi\'on de dependencias de terceros. En el \textit{Contexto de Sistema}, el artefacto de \textit{Democra} se conecta a sistemas externos, los cuales deben ser mapeados de forma expl\'icita:
\begin{itemize}
    \item \textbf{Motor de Mensajer\'ia Pub/Sub (WebSockets):} Un sistema externo (como \textit{Pusher} o \textit{Socket.io Redis}) utilizado para la propagaci\'on de eventos en tiempo real.
    \item \textbf{Servicio Transaccional de Env\'io de Correos:} Una API externa (ej. \textit{SendGrid} o \textit{AWS SES}) que gestiona la recuperaci\'on de contrase\~nas y el env\'io de notificaciones de auditor\'ia.
\end{itemize}
Esta documentaci\'on de dependencias satelitales es cr\'itica; dictamina que \textit{Democra} no es un monolito aislado, sino un orquestador que rutea peticiones de alta disponibilidad a trav\'es de internet.
